\documentclass[12pt, letterpaper]{article}

% - imports
\usepackage[utf8]{inputenc}
\usepackage{enumitem}
\usepackage[backend=biber]{biblatex-chicago}
\addbibresource{../bibliography.bib}
\usepackage{url}
\usepackage{parskip}
\usepackage{indentfirst}
\usepackage{parallel}

\usepackage[doublespacing]{setspace}

\title{\vspace*{-72pt} Working Title}
\author{Crystal Mandal}
\date{Last Edited \today}

\begin{document}

  \maketitle

  % \section*{Question 1}

  The image of the piano has had a long history of social signifier.
  The first pianos were invented in the Medici Court in the late 17th
  century \autocite*{piano-social-history} ; from the beginning, then, the piano has signified wealth.
  By the mid-1800's, the piano had shifted to a middle-class instrument in London. The novelty had
  worn off, and London's (aristocratic) musical scenes expected a certain degree of skill from the
  performer, leaving the playing of the piano to the professional musicians. As a consequence, the
  expectation that an aristocrat invite a professional - often international - performer to their
  salon slowly priced out much of the aristocracy from private musical
  affairs.\autocite*[61]{london-social-piano} Filling in this musical vacuum is the middle-class or
  bourgeois. Historian Richard Leppert argues that, ``[t]o the Victorian bourgeoisie the
  domestic [piano] was as essential as the dining room table".\autocite*[153]{sight-of-sound}
  Further, Leppert believes that, to the Victorian bourgeoisie, ``[the piano's] physical
  presence commonly fetishized materiality, ... [while] at the same time, the music to
be played on the instrument was valorized precisely because of its
immateriality".\autocite*[155]{sight-of-sound}
The piano is thus, by its absence, positioned as class signifier for the working classes -
that is, the working classes cannot afford the luxury of the piano.

Frederic Rzewski is by no means a conventional composer. He was born in 1938, and thus many
of his teachers were involved in the extreme modernist musics (and politics) of the early
20th century. His more notable teachers include American Composers Roger Sessions
and Milton Babbit, as well as notably politically impassioned Luigi Dallapiccolla. Rzewski's
musical language thus is heavily influenced by these 20th-century techniques, especially
serial techniques, modal writing, and poly\*- tonality. More important to Rzewski's impact
and sound is his engagement - particularly in the 60's and 70's - with the contemporary
leftist and revolutionary politics of the time. With titles, lyrics, and, most importantly,
musical thematic references, Rzewski places his music not in contemporary academic spheres -
with Stockhausen, or Cage - but rather with revolutionary spheres - with Pete Seeger, Sergio
Ortega, Hanns Eissler.

In the Early Seventies, Rzewski met and became familiar with Pete Seeger, who urged Rzewski
to perform in concert pieces that the audience would know or recognise. Seeger believed
it was essential that the audience have an avenue for participation - physically or mentally -
in concert music. \autocite*{rzewski-plays-rzewski} This advice then informed much of Rzewski's
composition. Many of his pieces make explicit reference to contemporary musics, particularly of
class struggle. In 1974, Frederic Rzewski and Ursula Oppens attended a concert given by Chilean
group Inti-Illimani at Hunter's College where they both heard Sergio Ortega's song
\textit{¡El pueblo unido jamás será vencido!} for the first time. Rzewski recalls walking out from the
performance still singing the melody, stating ``it never left us from that time
on". \autocite*{rzewski-plays-rzewski}
\clearpage{}

\nocite{*}
\printbibliography

\end{document}
